This dissertation presents the formulation and simulation of a Rendezvous and Docking or Berthing mission conducted by a robotic manipulator-type spacecraft, emphasizing the coherent sequencing of mission steps and the consistency between models, reference frames, and control laws throughout the operational profile. In this work, the PD control law was applied. The PD controller was implemented in the MATLAB/Simulink software environment. The controller itself is a well-established method in control theory and was not developed here; the contribution lies, in fact, in its implementation and application to the specific simulation scenario. The long-distance phase is represented by a Hohmann transfer between coplanar circular orbits, used to establish initial conditions for the short-distance approach; in this phase, the relative motion is described by the linearized Hill–Clohessy–Wiltshire equations and regulated by proportional-derivative control, while the pursuer's attitude is modeled by Euler's rigid body equation. In the final approach, the dynamic model is switched to a coupled spacecraft-manipulator system, obtained through Lagrange equations, and the PD control law acts to mitigate perturbations induced by the movement of the joints in the spacecraft's attitude. The robotic manipulator is controlled by the reference obtained through inverse kinematics. Performance is evaluated by the position error of the end effector, the coherence of the joint rotation, the actuation speed, and the maintenance of the pursuer's attitude and translation relative to the target. The results obtained indicate that the proposed architecture allows for phase transitions with stability and state coherence, enabling controlled approach and the required alignment in the studied scenario. The perturbation observed in the switching of the dynamic models in the final approach is noted. The limitations of the study include the use of linearized relative dynamics for coplanar circular orbits and the absence of environmental perturbations, which limits direct generalization to more complex operational scenarios.